\section{Discussion}
\label{sec:discussion}

\paragraph{Why narrative structure, and why does role assignment protect?}
The factorial results separate the ``storytelling'' bundle into a strong
structure effect (\Nfac{}) and a counter-intuitive protective effect (\Rfac{}).
A post-hoc account consistent with both is the elaboration-channel
hypothesis: \Nfac{}
increases elaboration (length, narrative density), which increases both the
probability and the perceived legitimacy of harmful continuations; \Rfac{}, by
contrast, constrains the model into an identity whose stance (an ``expert
advisor'') is at odds with producing harmful instructions in our protocol. The
replication of \Rfac's sign across models argues against a fluke, but its
sensitivity to the scoring caliber (it reverses under latent-label scoring)
demands caution; we report both.

\paragraph{Why do open-source guards fail on framing?}
ShieldGemma, WildGuard, and GradSafe all score below chance on framed inputs.
This is consistent with their training objective (harmfulness of direct
requests) being systematically mismatched with framing, which is
safe-looking. It motivates a \emph{knowledge} framing of the detection problem:
the relevant signal is not ``is this text harmful'' but ``does this structure
package harmful intent,'' which is exactly the signal our narrative and intent
branches measure.

\paragraph{Threshold detection versus mechanism-aware scoring.}
The benign-FPR failure (Section~\ref{sec:fusion-boundary}) is instructive rather
than incidental. It shows that a calibrated threshold is only meaningful inside the
calibration distribution, and that compliant benign behavior is statistically
confounded with framing signals. Read as a \emph{knowledge boundary} rather
than a defect, this is a substantive empirical finding about the framework's
applicability domain: the calibration pool (refusal responses) and the benign
pool (compliant responses) occupy different regions of the fused score space,
so the two distributions are not mutually calibrating. It is exactly this kind
of domain knowledge that our pre-registered defense gate G2 was designed to
elicit, and we report it in full. We argue that the honest contribution of a
fusion framework in this setting is a \emph{discriminative risk score with
auditable branch attributions} --- not a turnkey threshold --- and we present
the paper's detection claims accordingly. A direct path forward is
adaptive or online recalibration on deployment traffic, which we leave to
future work. This aligns with the journal's
orientation toward interpretable, decision-support knowledge systems
\citep{srd,realignment}.

\paragraph{Generalizability of the methodology.}
The contribution is best read as a \emph{methodology} rather than a tuned
detector. The four knowledge sources --- intent semantics, narrative structure,
acoustic prosody, and scoring uncertainty --- are defined from first principles
about the response and the scoring process, not about a specific model or
language; the causal protocol and the paired cross-modal measurement are
likewise model-agnostic. We therefore expect the framework to transfer to
other multimodal model audits (e.g., vision-language assistants whose safety
responses also admit packaging structure), to other high-resource languages,
and to high-stakes voice applications where an auditable, single-model risk
score is preferable to a stack of opaque classifiers. What transfers is the
\emph{structure}: mutually non-overlapping, auditable knowledge sources,
out-of-fold scoring, and honest boundary reporting. We have not re-run the
full protocol on these settings, and we state this explicitly rather than
claiming measured transfer.

\paragraph{Implications for knowledge-based decision systems.}
For the decision-support community, three design lessons follow. First, a
safety decision over a model response is a multi-source fusion problem. The
four signals we fuse are heterogeneous knowledge sources with different
reliability (the intent and narrative branches carry the signal; the acoustic
branch is ``not measured'', not ``ineffective''). Explicit calibration
plus an auditable fusion layer lets a system operator trace which source
drove a score. Second, the boundary between in-distribution discrimination and
threshold transfer is itself decision-relevant knowledge. The calibration pool
(adaptive refusals) and the benign pool (compliant responses) live in different
regions of the fused score space. A decision system must either keep its
operating context inside the calibrated region or re-calibrate online when the
deployment distribution drifts toward compliant benign traffic. Third, honest
boundary reporting is a decision-support feature, not a weakness: a score with
a documented failure region supports risk-informed routing (e.g., escalating
scores that fall outside the calibrated envelope to a human reviewer) better
than a threshold with an unmeasured false-positive rate. These are the
operational postures that follow from our measurement
\citep{guo2024survey,shafer1976,xu2025promethee}.
